%% kapa: exam=1 round=28 year=2017 period=1 form=A subject=civil_law subject_ko=민법 seq=24 number=24
%% source: Q-Net 감정평가사 기출문제 (cst00309) artlSeq=5199603 file=extracted/1008/5199603/2017년 제28회 감정평가사 1차시험 문제지/1교시 A형.pdf
%% conversion: kordoc 4.24.51 -> tools/gen1.py (auto); review=proofread-sonnet; flags=none
\begin{question}{24}
부동산 등기에 관한 설명으로 옳지 않은 것은? (다툼이 있으면 판례에 따름)
\begin{choices}
  \item 청구권보전을 위한 가등기에 기하여 본등기가 되더라도 본등기에 의한 물권변동의 효력이 가등기한 때로 소급하지 않는다.
  \item 소유권이전청구권이 정지조건부인 경우에도 가등기는 가능하다.
  \item 먼저 된 유효한 소유권보존등기로 인해 뒤에 된 이중보존등기가 무효인 경우, 뒤에 된 등기를 근거로 등기부취득시효를 주장할 수 있다.
  \item 등기되어 있는 3층 건물이 멸실되자 5층 건물을 신축하였으나 종전 등기를 그대로 사용하는 경우 이 등기는 무효이다.
  \item 등기명의인 표시변경등기는 권리변동을 가져오는 것은 아니다.
\end{choices}
\end{question}
\answer{3}
